% !TEX root = ../main.tex

\chapter{感知--控制闭环的自主任务执行}

\section{感知--控制闭环架构概述}

前两章分别解决了人形机器人全身控制中的数据来源问题和底层控制问题，而要让机器人在开放环境中完成真实任务，还需要将感知、规划与控制整合为闭环系统。本文面向动态交互和服务型任务，构建了统一的感知--控制闭环架构：感知模块负责识别环境状态与目标对象，任务规划模块负责将高层目标转化为结构化动作意图，控制模块则基于第 3 章提出的上下肢解耦控制框架完成具体动作执行，并将执行结果反馈给上层模块用于下一步决策。

该架构的关键不在于替代视觉感知或任务规划算法本身，而在于建立统一的接口规范。对于下肢控制器，上层模块输出目标方向、期望速度或局部导航目标；对于上肢控制器，上层模块输出末端目标位姿、关键轨迹路标点或抓取姿态约束。通过这种方式，来自不同任务的感知与规划结果可以在同一控制框架内复用，从而提高系统的通用性和工程效率。

\section{动态球拍运动控制任务}

\subsection{任务背景与系统架构}

CyboRacket 任务面向高速动态目标交互场景，要求机器人在较短时间内完成目标检测、落点预测、移动调整和挥拍动作生成。与静态抓取任务相比，该任务对系统时延、全身协调以及轨迹生成稳定性提出了更高要求：一方面，机器人需要根据高速运动目标的状态实时更新击球策略；另一方面，下肢必须在有限时间内调整站位，上肢则需要在满足可达性约束的同时完成挥拍动作。

在该任务中，本文关注的核心贡献在于全身控制接口与真机部署。视觉感知和轨迹预测模块提供球体状态估计与预期击球点信息，本文则负责将这些中间结果映射到第 3 章所定义的上下肢控制接口中，使机器人能够根据任务需要同时完成移动和挥拍。该过程体现了本文控制框架在动态任务中的可插拔性与闭环执行能力。

\subsection{视觉感知与轨迹预测接口}

为了支撑动态球拍任务，视觉系统需要以较高频率对目标进行检测，并通过物理模型或学习模型预测球体未来轨迹。对本文而言，这些感知结果并不直接作用于关节空间，而是被转换为对全身控制器更友好的中间表示，例如预期击球区域、建议机身站位以及击球时刻对应的上肢末端目标。

在下肢侧，预测轨迹被转换为局部导航指令，例如前向速度、横向位移或机身朝向调整需求，以驱动 RL locomotion 策略将机器人移动到合适击球区域；在上肢侧，预测击球点与拍面姿态要求被转化为末端位姿目标，并交由上肢轨迹生成和逆运动学模块求解。通过这种接口标准化方式，视觉模块无需直接考虑机器人全身动力学细节，而控制模块也无需依赖感知算法内部实现。

\subsection{全身运动生成与实验验证}

在获得感知输出后，系统首先确定机器人相对于目标轨迹的站位需求，并调用下肢控制器执行位置调整；随后，上肢控制器根据击球点、预期击球方向和时间窗口生成挥拍轨迹；最后，协调模块依据下肢当前支撑状态和上肢动作幅度进行姿态补偿，使机器人能够在移动与挥拍并行时保持整体稳定。

对于仿真验证，本文建议关注击球成功率、拦截位置误差和动作生成时延等指标，以衡量系统在动态任务中的闭环性能；对于真机验证，则应重点考察纯 onboard 感知条件下的持续执行能力、动作重复性和系统恢复能力。由于相关实验统计仍在整理中，本节先保留“目标速度大于 10\,m/s、系统总反应时间小于 500\,ms”等占位描述，后续可根据最终实验记录进行补全。

\section{语音指令服务任务执行}

\subsection{系统架构与任务流程}

除动态交互任务外，本文还面向具身服务机器人场景构建了语音指令服务任务执行系统。该任务以“给我拿瓶水”等自然语言请求为例，要求机器人从高层语义输入出发，完成目标理解、路径规划、目标感知、抓取执行和物品递送等完整流程。与 CyboRacket 相比，该任务对语义理解与多阶段流程组织要求更高，但对全身控制框架的复用同样具有代表性。

在系统实现中，语音输入首先通过飞书 API 或其他交互接口转换为文本指令，高层规划模块基于 OpenClaw 将用户意图拆解为导航、检测、抓取和递送等结构化子任务；随后，SAM3 等视觉模块负责输出场景分割结果和目标位姿估计；最后，这些任务结果被映射到第 3 章统一控制接口中，以驱动机器人逐步完成任务执行。

\subsection{任务规划与视觉感知}

任务规划模块的作用是将高层语义目标转换为控制系统可执行的阶段性任务。例如，“拿瓶水”这一请求可以被分解为接近目标区域、搜索瓶体、定位抓取位姿、抓取目标、移动至交付位置和递送目标等多个阶段。规划器输出的每个阶段既包含动作类别，也包含必要的空间目标与阶段切换条件。

视觉感知模块则负责在各任务阶段提供实时环境信息。在导航阶段，需要识别目标物体或目标区域在环境中的位置关系；在操作阶段，需要输出物体的几何位姿、抓取接触区域和周围障碍物分布；在递送阶段，还需结合人机交互状态确定最终交付姿态。通过将这些信息统一组织为任务阶段约束和末端位姿目标，系统可以顺畅接入本文提出的上、下肢控制框架。

\subsection{全身控制集成与执行}

在服务任务执行过程中，下肢控制器负责依据目标方位和路径规划结果生成基础移动行为。对于较简单的近距离任务，可将目标方向直接映射为前进与旋转速度命令；对于需要更复杂避障或局部修正的场景，则可通过局部导航模块先生成中间路径点，再由 RL locomotion 策略完成连续移动。

上肢控制器则根据目标物体位姿和抓取计划生成抓取轨迹，并在抓取后切换到持物递送模式。为保证移动抓取场景下的稳定性，系统会结合第 3 章中的全身协调机制，对机身姿态和重心变化进行补偿。当机器人一边移动一边进行视觉调整或上肢预摆时，该机制能够降低大幅动作对行走稳定性的影响。

\subsection{实验验证}

服务任务的实验验证可从任务成功率、平均执行时间、阶段失败原因和人机交互体验等角度展开。对于室内办公环境和 cluttered 桌面等典型场景，系统应重点分析视觉识别失效、抓取失败、移动误差积累以及阶段切换不及时等问题。通过将失败案例划分为感知失败、规划失败和控制失败三类，可以更清晰地定位系统性能瓶颈。

在当前论文初稿中，相关实验数据尚未最终整理，因此本文保留任务成功率、平均执行时长与失败占比等指标占位。后续完成系统测试后，可进一步在本节补充表格和案例图示，并据此分析当前系统在语义理解、复杂环境适应性和长时序任务连续执行方面的局限。

\section{本章小结}

本章围绕两个代表性任务展示了本文方法在感知--控制闭环场景中的应用方式。CyboRacket 任务验证了本文框架在动态目标、高时效要求场景中的全身控制能力，OpenClaw 服务任务则体现了统一控制接口在语义任务执行链路中的适配性。两类任务共同说明，本文提出的技术路线并非局限于离线动作重现，而是能够支撑人形机器人从感知结果到自主动作执行的系统级闭环。
